\begin{lemma}
\label{editorial-source-lemma-integral-no-inclusion}
Suppose $R \to S$ is integral.
Let $\mathfrak q, \mathfrak q' \in \Spec(S)$
be distinct primes
having the same image in $\Spec(R)$.
Then neither $\mathfrak q \subset \mathfrak q'$
nor $\mathfrak q' \subset \mathfrak q$.
\end{lemma}

\begin{proof}
Let $\mathfrak p=\varphi^{-1}(\mathfrak q)=
\varphi^{-1}(\mathfrak q')$ and set $U=R\setminus\mathfrak p$.
Retain the fibre algebra
$A=S\otimes_R\kappa(\mathfrak p)\cong U^{-1}S/\mathfrak pU^{-1}S$.
The two ideals are sent to
$U^{-1}\mathfrak q/\mathfrak pU^{-1}S$ and
$U^{-1}\mathfrak q'/\mathfrak pU^{-1}S$. Localization and quotient
preserve inclusion, and contraction recovers each original prime:
$s/1\in U^{-1}\mathfrak q$ means
$\varphi(u)s\in\mathfrak q$ for some $u\in U$, which by primality
and $\varphi(u)\notin\mathfrak q$ is equivalent to $s\in\mathfrak q$.
Thus the two resulting primes are distinct. Lemma
\ref{editorial-source-lemma-base-change-integral} makes $A$ integral over
$\kappa(\mathfrak p)$, so Lemma \ref{editorial-source-lemma-integral-over-field}
makes both primes maximal. An inclusion between them would force
equality, and contraction would contradict the original distinction.
\end{proof}